% String recovery examples figure using actual chunked decompiled output.
% Compile standalone or \input{} into a larger document.

\documentclass[border=10pt]{standalone}
\usepackage[T1]{fontenc}
\usepackage{lmodern}
\usepackage{listings}
\usepackage{xcolor}
\usepackage{tcolorbox}
\tcbuselibrary{listings, skins, breakable}
\usepackage{varwidth}

% Colors
\definecolor{codebg}{HTML}{F7F7F7}
\definecolor{commentgreen}{HTML}{3D7B3D}
\definecolor{keywordblue}{HTML}{0033B3}
\definecolor{stringbrown}{HTML}{8B4513}
\definecolor{resolved}{HTML}{1A6B1A}
\definecolor{phase1col}{HTML}{3366AA}
\definecolor{phase2col}{HTML}{AA6633}
\definecolor{phase3col}{HTML}{7733AA}
\definecolor{phase4col}{HTML}{AA3333}
\definecolor{labelgray}{HTML}{666666}

% Listing style for C pseudocode
\lstdefinestyle{cstyle}{
    language=C,
    basicstyle=\ttfamily\fontsize{7.5}{9.5}\selectfont,
    keywordstyle=\color{keywordblue},
    commentstyle=\color{commentgreen},
    stringstyle=\color{stringbrown},
    showstringspaces=false,
    breaklines=true,
    columns=flexible,
    tabsize=2,
    xleftmargin=2pt,
    xrightmargin=2pt,
    aboveskip=0pt,
    belowskip=0pt,
    escapeinside={(@}{@)},
    morekeywords={undefined8, undefined1, ulong, uint8},
}

% Listing style for annotations
\lstdefinestyle{annot}{
    basicstyle=\ttfamily\fontsize{7.5}{9.5}\selectfont\color{commentgreen},
    showstringspaces=false,
    breaklines=true,
    columns=flexible,
    xleftmargin=2pt,
    xrightmargin=2pt,
    aboveskip=0pt,
    belowskip=0pt,
}

% Uniform gap between phases
\newcommand{\phasegap}{\vspace{8pt}}

\begin{document}

\begin{varwidth}{16cm}

% Phase 1
\begin{tcolorbox}[
    title={\sffamily\bfseries Heuristic 1: Defined Data Scan},
    colframe=phase1col, colback=codebg, colbacktitle=phase1col!15!white,
    coltitle=phase1col, fonttitle=\small,
    width=16cm, boxrule=0.5pt, left=3pt, right=3pt, top=2pt, bottom=2pt,
    subtitle style={colback=phase1col!8!white, coltitle=phase1col},
]
\small\sffamily Scans Ghidra's typed data items once per binary.
\texttt{s\_*} labels yield strings directly; \texttt{PTR\_s\_*} labels are pointer-dereferenced.

\tcblower

{\scriptsize\sffamily\color{labelgray} From \texttt{tailscale/tailscale} --- \texttt{client/systray.setAppIcon} (stripped build):}
\vspace{2pt}
\begin{lstlisting}[style=annot]
// Strings:
//   (@\textcolor{phase1col}{\textbf{s\_loading\_00f7a900}}@) = "loading"
\end{lstlisting}
\vspace{2pt}
{\scriptsize\sffamily\color{labelgray} C pseudocode --- recovered string used in a comparison:}
\begin{lstlisting}[style=cstyle]
if ((in_stack_00000008 == (@\textcolor{phase1col}{\textbf{s\_loading\_00f7a900}}@))
    && (DAT_00f7a908 == cStack0000000000000010)) {
  tailscale_com_client_systray_startLoadingAnimation();
}
\end{lstlisting}
{\scriptsize\sffamily\color{labelgray} Ghidra typed \texttt{s\_loading} as string data at \texttt{0xf7a900}; the pipeline reads \texttt{"loading"} directly from that address.}
\end{tcolorbox}

\phasegap

% Phase 2
\begin{tcolorbox}[
    title={\sffamily\bfseries Heuristic 2: Undefined Data Resolution},
    colframe=phase2col, colback=codebg, colbacktitle=phase2col!15!white,
    coltitle=phase2col, fonttitle=\small,
    width=16cm, boxrule=0.5pt, left=3pt, right=3pt, top=2pt, bottom=2pt,
    subtitle style={colback=phase2col!8!white, coltitle=phase2col},
]
\small\sffamily Resolves \texttt{PTR\_s\_*} / \texttt{s\_*} references in C pseudocode absent from the data listing.
Parses hex suffix from name; reads Go \texttt{(ptr, len)} pairs for precise, length-delimited extraction.

\tcblower

{\scriptsize\sffamily\color{labelgray} From \texttt{usememos/memos} --- \texttt{plugin/filter.NewAttachmentSchema} (default build):}
\vspace{2pt}
\begin{lstlisting}[style=annot]
// Strings:
//   (@\textcolor{phase2col}{\textbf{PTR\_s\_attachmentindex\_html\_..\_01aad170}}@) = ["attachment", "memo_id"]
//   (@\textcolor{phase2col}{\textbf{PTR\_s\_attachmentindex\_html\_..\_01aacfc0}}@) = ["attachment", "filename",
//     "mime_type", "scalar", "string", "attachment", "type",
//     "create_time", "scalar", "timestamp", ...]
\end{lstlisting}
\vspace{2pt}
{\scriptsize\sffamily\color{labelgray} C pseudocode referencing the symbols (initializing map entries):}
\begin{lstlisting}[style=cstyle]
local_68 = (@\textcolor{phase2col}{\textbf{PTR\_s\_attachmentindex\_html\_..\_01aacfc0}}@);
uStack_60 = _UNK_01aacfc8;
  ...
runtime_mapassign(in_RDI, in_RSI, key, &PTR_DAT_01a76720);
\end{lstlisting}
{\scriptsize\sffamily\color{labelgray} Hex suffix \texttt{\_01aacfc0} parsed $\to$ memory at \texttt{0x1aacfc0} read as Go \texttt{(ptr, len)} pairs $\to$ 17 domain-specific strings recovered.}
\end{tcolorbox}

\phasegap

% Phase 3
\begin{tcolorbox}[
    title={\sffamily\bfseries Heuristic 3: DAT Symbols},
    colframe=phase3col, colback=codebg, colbacktitle=phase3col!15!white,
    coltitle=phase3col, fonttitle=\small,
    width=16cm, boxrule=0.5pt, left=3pt, right=3pt, top=2pt, bottom=2pt,
    subtitle style={colback=phase3col!8!white, coltitle=phase3col},
]
\small\sffamily Resolves \texttt{DAT\_*} labels (unclassified data).
Detects paired string lengths from adjacent stack locals or function arguments.

\tcblower

{\scriptsize\sffamily\color{labelgray} From \texttt{ollama/ollama} --- \texttt{main.main.func1} in \texttt{pull-progress} (default build):}
\vspace{2pt}
\begin{lstlisting}[style=annot]
// Strings:
//   (@\textcolor{phase3col}{\textbf{DAT\_007ac68b}}@) = "Progress: status=%v, total=%v, completed=%v"
\end{lstlisting}
\vspace{2pt}
{\scriptsize\sffamily\color{labelgray} C pseudocode --- used as a format string argument:}
\begin{lstlisting}[style=cstyle]
w.tab = (internal_abi_ITab *)0x2c;       (@\textcolor{phase3col}{// length = 44}@)
format.len = (int)&(@\textcolor{phase3col}{\textbf{DAT\_007ac68b}}@);
format.str = extraout_RDX;
fmt_Fprintf(w, format, in_stack, 3, ...);
\end{lstlisting}
{\scriptsize\sffamily\color{labelgray} Hex address parsed from label. Length \texttt{0x2c} (44) detected from adjacent assignment, matching the 44-byte format string.}
\end{tcolorbox}

\phasegap

% Phase 4
\begin{tcolorbox}[
    title={\sffamily\bfseries Phase 4: Hex Literals},
    colframe=phase4col, colback=codebg, colbacktitle=phase4col!15!white,
    coltitle=phase4col, fonttitle=\small,
    width=16cm, boxrule=0.5pt, left=3pt, right=3pt, top=2pt, bottom=2pt,
    subtitle style={colback=phase4col!8!white, coltitle=phase4col},
]
\small\sffamily Resolves bare hex constants in C pseudocode that fall within \texttt{.rodata} address ranges.

\tcblower

{\scriptsize\sffamily\color{labelgray} From \texttt{argoproj/argo-workflows} --- \texttt{v1alpha1.ClusterWorkflowTemplate.GetResourceScope} (stripped build):}
\vspace{2pt}
\begin{lstlisting}[style=annot]
// Strings:
//   (@\textcolor{phase4col}{\textbf{0x3171032}}@) = "cluster"
\end{lstlisting}
\vspace{2pt}
{\scriptsize\sffamily\color{labelgray} C pseudocode --- bare hex address returned as a string pointer:}
\begin{lstlisting}[style=cstyle]
undefined8 ...GetResourceScope(void) {
  return (@\textcolor{phase4col}{\textbf{0x3171032}}@);
}
\end{lstlisting}
{\scriptsize\sffamily\color{labelgray} Address \texttt{0x3171032} falls within \texttt{.rodata}. The decompiler shows only a numeric constant; Phase~4 resolves it to \texttt{"cluster"}.}
\end{tcolorbox}

\end{varwidth}

\end{document}
